%-----------------------------------------------
% DOCUMENT PREAMBLE
%-----------------------------------------------

\documentclass[12pt]{article}

%call general packages for use here
\usepackage[utf8]{inputenc}
\usepackage[T1]{fontenc}
\usepackage{mathptmx}
\usepackage{amsmath}
\usepackage{amssymb}

%in-document hyperlink control
\usepackage{hyperref}

%inserting image files
%\usepackage{graphicx}
%\graphicspath{{images/}}

%margin control
\usepackage{geometry}
\geometry{
  inner=1.25in,
  outer=1.25in,
  top=1.25in,
  bottom=1.25in,
  bindingoffset=.5cm,
  %showframe %uncomment to see the type block
}

%section format control
\usepackage{titlesec}
\newcommand{\sectionbreak}{\clearpage}
\setlength{\parindent}{0pt}
%set paragraph indentation

%\usepackage[doublespacing]{setspace} %set doublespacing

%improves typeset quality
\usepackage{microtype}

%-----------------------------------------------
% BEGIN DOCUMENT
%-----------------------------------------------

\begin{document}

%-----------------------------------------------
% TITLE PAGE
%-----------------------------------------------

\thispagestyle{empty}
\begin{center}
  \vspace*{1.5in}
  THESIS PROPOSAL\\[0.25in]
  By\\[0.5in]
  Asher Lucas-Cuddeback\\[0.25in]
  00398526\\[1.75in]
  Submitted to Northern Michigan University to propose\\
  a course of thesis work pursuant to the degree of\\[0.5in]
  MASTER OF SCIENCE IN MATHEMATICS\\[1.5in]
  Date: 08/25/2026
\end{center}

%-----------------------------------------------
% THESIS PROPOSAL CONTENT
%-----------------------------------------------

\newpage
\thispagestyle{empty}

\paragraph{} \textbf{Motivation.} Quadratic reciprocity is one of the
crown jewels of elementary number theory, but its higher-degree
analogues -- cubic, quartic, and beyond -- require genuinely new
machinery to even state cleanly, let alone prove. Gauss himself
recognized that cubic reciprocity could not be phrased purely in
terms of ordinary integers, and instead requires working in the ring
of Eisenstein integers $\mathbb{Z}[\omega]$, where $\omega$ is a
primitive cube root of unity. This thesis aims to understand this
machinery through the modern lens of Galois and Class Field theory.
The main drivers will be providing specific, concrete understandings
of the concepts at hand.

\paragraph{} \textbf{Background.} The cyclotomic field
$K = \mathbb{Q}(\omega)$ has ring of integers $\mathbb{Z}[\omega]$,
a unique factorization domain in which rational primes $p \equiv 1
\pmod 3$ split, while $p \equiv 2 \pmod 3$ remain inert. For a prime
$\pi \in \mathbb{Z}[\omega]$ of norm $p$, the cubic residue character
$\left(\frac{\alpha}{\pi}\right)_3$ generalizes the Legendre symbol,
taking values in $\{1, \omega, \omega^2\}$ according to whether
$\alpha$ is a cubic residue mod $\pi$. The reciprocity law relates
$\left(\frac{\alpha}{\pi}\right)_3$ to $\left(\frac{\pi}{\alpha}\right)_3$
up to a unit correction, mirroring the classical quadratic case but
requiring the extra structure of the Galois group
$\mathrm{Gal}(K/\mathbb{Q}) \cong \mathbb{Z}/2\mathbb{Z}$ and its
interaction with decomposition and inertia groups in ramified
extensions.

\paragraph{} \textbf{Solvability.} A Galois theoretic treatment of
reciprocity at degrees \(n\geq 3\) uncovers that the traditional statement
of quadratic reciprocity as congruences modulo an integer is the exceptional
case. To state cubic and quartic (biquadratic in the language of
Gauss) reciprocity
in the case of a splitting problem for polynomials with
\textit{solvable} Galois groups,
we can not rely solely on congruences. For example, the polynomial
\(f(x)=x^3-2\) splits modulo \(p\) if and only if \(p=a^2+27b^2\) for
integer \(a,b\). The failure of the congruence understanding lies with the
increasing complexity of the underlying Galois group. The function
\(f\) has splitting field \( \mathbb{Q}(\sqrt[3]{2}, \omega) \). When
we compute \( \mathrm{Gal}(\mathbb{Q}(\sqrt[3]{2}, \omega) /
\mathbb{Q}) \) we find

\[ \mathrm{Gal}(\mathbb{Q}(\sqrt[3]{2}, \omega) / \mathbb{Q}) \cong S_3\]

A congruence condition arises from Galois extensions which are
cyclic. So if we look closer, we
can see their ghost.
Let \(H = A_3 \trianglelefteq S_3\), the unique subgroup of order 3.
Its fixed field is \(\mathbb{Q}(\omega)\), and since \(H\) is normal,
\(\mathbb{Q}(\omega)/\mathbb{Q}\) is itself Galois, with
\[ \mathrm{Gal}(\mathbb{Q}(\omega)/\mathbb{Q}) \cong S_3/A_3 \cong
\mathbb{Z}/2\mathbb{Z}. \]
This cyclic quotient is the ghost of a congruence condition: it is
precisely what forces \(p \equiv 1 \pmod 3\) whenever \(f\) splits
mod \(p\), matching the necessary condition noted in
Example~1.1.2. But this quotient only records whether \(p\) splits in
\(\mathbb{Q}(\omega)\); it is silent on the finer question of whether
\(\sqrt[3]{2}\) itself becomes rational mod \(p\). That information
lives in the non-normal subgroups of order 2 in \(S_3\), whose fixed
fields \(\mathbb{Q}(\sqrt[3]{2})\), \(\mathbb{Q}(\omega\sqrt[3]{2})\),
\(\mathbb{Q}(\omega^2\sqrt[3]{2})\) are conjugate to one another but
not normal over \(\mathbb{Q}\). Because no single one of them is
Galois over \(\mathbb{Q}\), no congruence on \(p\) alone can detect
which case occurs; the obstruction must instead be read off inside
\(\mathbb{Z}[\omega]\), from whether a prime above \(p\) admits a
generator that is a perfect cube. It is exactly this refinement --
invisible at the level of \(p \bmod 3\), but visible in the
arithmetic of \(\mathbb{Z}[\omega]\) -- that manifests as the
condition \(p = a^2+27b^2\) in place of a congruence.
\vfill

\paragraph{} REFERENCES

\paragraph{} Cox, David A. \textit{Primes of the Form $x^2 + ny^2$:
Fermat, Class Field Theory, and Complex Multiplication}. Wiley-
Interscience, 1989.
\paragraph{} Weinstein, Jared. "Reciprocity Laws and Galois
Representations: Recent Breakthroughs." \textit{Bull. Amer. Math.
Soc.} 53 (2016), 1--39.

%-----------------------------------------------
% END DOCUMENT
%-----------------------------------------------

\end{document}
